% --- Slide 1: SFT train ---
\begin{frame}[t]{Dữ liệu huấn luyện SFT (500 mẫu)}
  \begin{columns}[T,onlytextwidth]
    \begin{column}{0.48\textwidth}
      \begin{block}{Input --- \texttt{seed}}
        \scriptsize
        Thương hiệu: MiniGo Mart -- siêu thị mini, 8 cửa hàng HN\\
        Sản phẩm: Combo bữa tối 5 món cho 2 người\\
        Giá: 129.000đ (gốc 159.000đ, $-$19\%), freeship 2km\\
        Điều kiện: 17h--20h, đặt FB/Zalo, giới hạn 40 combo/ngày\\
        Đối tượng: 24--38 tuổi, nhân viên VP, Cầu Giấy -- HBT\\
        KPI: 50 inbox khảo sát, CTR $\geq$ 2,5\%
      \end{block}
    \end{column}
    \begin{column}{0.50\textwidth}
      \begin{block}{Output --- \texttt{content} (đã kiểm duyệt)}
        \scriptsize
        ``Tan làm rồi mà vẫn phải nghĩ \textit{tối nay ăn gì} là
        chuyện rất quen với dân văn phòng ở Hà Nội.\\[0.3em]
        MiniGo Mart đang thử nghiệm Combo bữa tối 5 món cho 2
        người, kiểu mua nhanh -- đủ món -- không phải đi xa
        [\ldots]\\[0.3em]
        Inbox ngay nếu bạn muốn trải nghiệm thử nhé.''
      \end{block}
    \end{column}
  \end{columns}
  \vspace{0.5em}
  \begin{block}{Cấu trúc mỗi mẫu}
    \scriptsize
    \texttt{instruction} (vai trò copywriter + ràng buộc phong cách)
    $+$ \texttt{title} $+$ \texttt{seed} $\xrightarrow{\text{học}}$ \texttt{content}
  \end{block}
\end{frame}

% --- Slide 2: RL / Dining Review ---
\begin{frame}[t]{Dữ liệu huấn luyện RL --- Dining Review (500 bài)}
  \begin{columns}[T,onlytextwidth]
    \begin{column}{0.60\textwidth}
      \begin{block}{Ví dụ bài viết thực}
        \scriptsize
        ``Địa chỉ trà chiều mới tại HN. Ks mới khai trương,
        trà chiều ngon hơn đứt các chỗ tại VN đã từng thử.
        Bánh ngon thật sự, scone vỏ giòn nhẹ, ruột mềm ẩm.
        Clotted cream béo và ngậy, chắc tự làm mới được thế.
        Ngồi dễ chịu, phục vụ dễ thương
        \#afternoontea \#hanoi \#diningreview''
      \end{block}
    \end{column}
    \begin{column}{0.37\textwidth}
      \begin{block}{Chỉ số tương tác thực}
        \scriptsize
        Likes: \textbf{125} \quad Love: \textbf{11}\\
        Comments: \textbf{50}\\
        Shares: \textbf{10}\\[0.4em]
        $e_{\text{raw}} = L + 0{,}5C + 1{,}5S + \ldots$\\
        $e = \ln(1 + e_{\text{raw}}) = \mathbf{5{,}68}$
      \end{block}
    \end{column}
  \end{columns}
  \vspace{0.5em}
  \begin{block}{Chiến lược Prefix Completion}
    \scriptsize
    Dataset \textbf{không có} cặp prompt--response rõ ràng.
    Giải pháp: \textbf{câu đầu tiên} của mỗi bài = \emph{prompt};
    phần còn lại do model sinh ra và hàm reward rule-based chấm điểm.
    Model học hoàn thiện bài post sao cho các tín hiệu hình thức
    tương quan với engagement cao trong dataset.
  \end{block}
\end{frame}

% --- Slide 3: Test set ---
\begin{frame}[t]{Tập kiểm thử chung (100 mẫu)}
  \begin{columns}[T,onlytextwidth]
    \begin{column}{0.55\textwidth}
      \begin{block}{Input --- \texttt{seed} (ví dụ thực)}
        \scriptsize
        Thương hiệu: Matcha Mây -- trà trái cây/matcha, 8 điểm bán TP.HCM\\
        Sản phẩm: Mid-Autumn Duo -- 2 ly 700ml + hộp quà mini\\
        Giá: 119.000đ (gốc 149.000đ, $-$20\%), freeship 3km\\
        Điều kiện: 15/8--15/9, GrabFood/inbox, 80 set/ngày\\
        Đối tượng: 22--35 tuổi, dân VP, Q1, Q3, PN, BT\\
        KPI: CTR $\geq$ 3,5\%, ROAS $\geq$ 4, GMV $\geq$ 180 triệu
      \end{block}
    \end{column}
    \begin{column}{0.42\textwidth}
      \begin{block}{\texttt{content} --- để trống}
        \scriptsize
        Mô hình tự sinh đầu ra từ seed.\\[0.4em]
        Pipeline chấm điểm 3 lớp:\\
        \textbf{F}aithfulness / \textbf{E}xpansion / \textbf{V}ibe\\[0.4em]
        Judge: \textbf{GPT-5.4 Azure}\\[0.4em]
        Dùng chung để so sánh\\
        tất cả mô hình (Baseline,\\
        SFT, RL, GPT-5.4-mini)
      \end{block}
    \end{column}
  \end{columns}
  \vspace{0.5em}
\end{frame}
